\section{Evaluación y evaluación de riesgos}

\subsection{Evaluación a nivel de modelo vs. evaluación a nivel de sistema}

La evaluación a nivel de modelo se centra en el comportamiento de un modelo aislado ante entradas específicas (por ejemplo, MMLU, pruebas de referencia de matemáticas). Los sistemas híbridos Agentic requieren \textbf{una evaluación del comportamiento del sistema de extremo a extremo}, cuya complejidad es muchísimo mayor que la de la evaluación a nivel de modelo.

El trabajo relacionado del equipo de Dawn Song incluye:
\begin{itemize}
    \item \textbf{DecodingTrust}: el primer marco de evaluación integral de la confiabilidad de los LLM (premio al mejor artículo de NeurIPS, premio al mejor artículo del año de la NSA)
    \item \textbf{MMDT}: marco de evaluación de la confiabilidad y la seguridad de los modelos fundacionales multimodales (ICLR 2025)
    \item \textbf{Evaluación de riesgos de Coding Agent} (NeurIPS 2024)
\end{itemize}

\subsection{AgentExploit: pruebas de equipo rojo de extremo a extremo para Agent de caja negra}

\begin{importantbox}{Marco AgentExploit}
En un entorno de caja negra, el atacante solo puede controlar las fuentes de datos externas (como páginas web o archivos), no puede modificar la consulta del usuario ni acceder a la estructura interna del Agent, y solo puede obtener retroalimentación binaria (éxito o fracaso del ataque).

El marco utiliza el método de \textbf{pruebas de fuzzing} (Fuzzing):
\begin{enumerate}
    \item Se parte de un conjunto inicial de instrucciones de ataque semilla
    \item Se mutan las semillas para generar nuevos ataques
    \item Se inyecta contenido malicioso en las fuentes de datos externas
    \item Se obtiene retroalimentación y se optimiza la selección de semillas y la estrategia de mutación
\end{enumerate}
Resultados experimentales: la tasa de éxito del ataque es el doble de la línea base manual, y los ataques generados tienen transferibilidad.
\end{importantbox}
